
\documentclass[a4paper,12pt,fleqn]{article}

% --- Core Packages ---
\usepackage[margin=2cm]{geometry}
\usepackage[utf8]{inputenc}
\usepackage[T1]{fontenc}
\usepackage{libertine} % Font matching Linux Libertine
\usepackage{amsmath, amssymb, amsthm}

% --- Inference Rules Package ---
\usepackage{bussproofs}

% --- Left-align display math ---
\setlength{\mathindent}{0pt}

% --- No paragraph indentation ---
\setlength{\parindent}{0pt}

% --- Custom Commands ---
\newcommand{\tack}{\vdash}
\newcommand{\tackl}{\dashv}

\begin{document}

\section{Coherence Semantics}
This document describes the semantics on which the Coherence language is based on.

\section{Basic cut-down version of the language}
This cutdown version does not have the full features of the language. It is described as below.

\subsection{Formal Syntax}

\subsubsection{Types}
\[
\begin{array}{lcll}
T & ::= & B \mid B\ \kappa & \text{(full coherence type)} \\[0.5em]
B & ::= & \mathtt{unit} \mid \mathtt{int} \mid \mathtt{bool} \mid \mathtt{Named}(N) \mid \mathtt{Actor}(A) & \text{(basic type)} \\[0.5em]
\kappa & ::= & \mathtt{ref} \mid \mathtt{val} \mid \mathtt{iso} \mid \mathtt{locked}\langle L \rangle & \text{(reference capabilities)} \\[0.5em]
\hat{T} & ::= & B  \mid \mathtt{struct}\ \{\ \overline{f : B;}\ \} & \text{(nameable type)}
\end{array}
\]
\\
TODO: Add explanation of this stuff
\\
\subsubsection{Top-level items}
These are items that occur in the top-level of a program. They include type definitions, top-level functions, and actor definitions.
\[
\begin{array}{lcll}
P & ::= & \overline{D} & \text{(program)} \\[0.5em]
D & ::= & \mathtt{type}\ N = \hat{T} & \text{(type definition)} \\
  & \mid & \mathtt{func}\ m(\overline{x : T}) \Rightarrow T\ \{\ \overline{s}\ \} & \text{(function definition)} \\
  & \mid & \mathtt{actor}\ A\ \{\ \overline{f : T};\ \overline{M}\ \} & \text{(actor definition)} \\[0.5em]
M & ::= & \mathtt{new}\ k(\overline{T\ \mathit{arg}})\ \{\ \overline{s}\ \} & \text{(constructor)} \\
  & \mid & \mathtt{func}\ m(\overline{T\ \mathit{arg}}) \Rightarrow T\ \{\ \overline{s}\ \} & \text{(method)} \\
  & \mid & \mathtt{be}\ b(\overline{T\ \mathit{arg}})\ \{\ \overline{s}\ \} & \text{(behaviour)}
\end{array}
\]

\subsubsection{Statements}
Statements are parts of the body that are terminated with a semi-colon in the body of a callable (function, behaviour, or constructor). They can have side effects and do not return a value.

\[
\begin{array}{lcll}
s & ::= & \mathtt{var}\ x : T = e\ ; & \text{(local declaration)} \\
  & \mid & f := e\ ; & \text{(member initialisation)} \\
  & \mid & e \to b(\overline{e})\ ; & \text{(behaviour call)} \\
  & \mid & \mathtt{OUT}\ e\ ; & \text{(print statement)} \\
  & \mid & \mathtt{if}\ (e)\ \{\ \overline{s}\ \}\ \mathtt{else}\ \{\ \overline{s}\ \} & \text{(if-else statement)} \\
  & \mid & \mathtt{while}\ (e)\ \{\ \overline{s}\ \} & \text{(while statement)} \\
  & \mid & \mathtt{atomic}\ \{\ \overline{s}\ \} & \text{(atomic section)}\\
  & \mid & \mathtt{return}\ e\ ; & \text{(return)}\\
  & \mid & e\ ; & \text{(expression whose value is discarded)}
\end{array}
\]

\subsubsection{Expressions}
These are parts of the program that evaluate to a value.

\[
\begin{array}{lcll}
e & ::= & () & \text{(unit)} \\
  & \mid & i & \text{(integer literal)} \\
  & \mid & \mathtt{true} \mid \mathtt{false} & \text{(boolean literal)} \\
  & \mid & x & \text{(variable)} \\
  & \mid & \{\ \overline{f = e;}\ \} : N & \text{(struct literal)} \\
  & \mid & \mathtt{new}\ \kappa\ [e]\ B\ (e) & \text{(allocation)} \\
  & \mid & \mathtt{new}\ A.k(\overline{e}) & \text{(actor construction)} \\
  & \mid & e[e] & \text{(pointer access)} \\
  & \mid & e.f & \text{(field access)} \\
  & \mid & e = e & \text{(assignment)} \\
  & \mid & m(\overline{e}) & \text{(function call)} \\
  & \mid & e\ \mathit{op}\ e & \text{(binary operation)}
\end{array}
\]

\section{Type checking}
This part of the document describes the type-checking rules carried out on the cut-down version of the language. In the formal rules below, assume that renaming has been carried out and so all names are distinct and conflict-free.

\subsection{Typing Environment}
The environment $\Gamma$ consists of the following:
\begin{itemize}
    \item $\Gamma_V$: a mapping from variable names to their types
    \item $\Gamma_L$: a boolean about whether the type-checker is in an atomic section
    \item $\Gamma_A$: a mapping from actor names to an object containing their information (constructors, functions and behaviours).
    \item $\Gamma_T$: mapping to named types.
    \item $\Gamma_F$: mapping from function names to their information.
    \item $\Gamma_{\text{this}}$: the type of the current actor.
\end{itemize}

\subsection{Expression typing}
I will not be focusing on binary operation typing as it is kind of simple and not worth mentioning.
\\
\\
\begin{tabular}{c@{\qquad}c@{\qquad}c}
    \AxiomC{}
    \RightLabel{(Int)}
    \UnaryInfC{$\Gamma \tack i : \text{int}$}
    \DisplayProof
&
    \AxiomC{}
    \RightLabel{(True)}
    \UnaryInfC{$\Gamma \tack \text{true} : \text{bool}$}
    \DisplayProof
&
    \AxiomC{}
    \RightLabel{(False)}
    \UnaryInfC{$\Gamma \tack \text{false} : \text{bool}$}
    \DisplayProof
\end{tabular}

\vspace{1em}

\begin{prooftree}
    \AxiomC{$\text{variable\_type}(x, \Gamma) = T$}
    \RightLabel{(Var)}
    \UnaryInfC{$\Gamma \tack x : T$}
\end{prooftree}

\vspace{1em}
\begin{prooftree}
    \AxiomC{$\text{get\_struct\_type}(N, \Gamma) = \text{struct } \{ \overline{f:B} \}$}
    \AxiomC{$\Gamma \tack \overline{e} : \overline{B'}$}
    \AxiomC{$\Gamma \tack \text{assignable}(\overline{B}, \overline{B'})$}
    \RightLabel{(Struct)}
    \TrinaryInfC{$\Gamma \tack (\{ \overline{f = e} \} : N) : N$}
\end{prooftree}

\begin{prooftree}
    \AxiomC{$\Gamma \tack e_s : \text{Int}$}
    \AxiomC{$\Gamma \tack e_d : B'$}
    \AxiomC{$\Gamma \tack \text{assignable}(B, B')$}
    \RightLabel{(Alloc)}
    \TrinaryInfC{$\Gamma \tack \text{new } \kappa[e_s] B(e_d) : \text{unalias}(B \ \kappa)$}
\end{prooftree}

\begin{prooftree}
    \AxiomC{$A \in \Gamma_A$}
    \AxiomC{$k \in A_{\text{ctor}}$}
    \AxiomC{$\Gamma \tack \overline{e} : \overline{T_{\text{arg}}}$}
    \AxiomC{$\Gamma \tack \text{type\_assignable}(k_{\text{param}}, \overline{T_{\text{arg}}})$}
    \RightLabel{(Actor creation)}
    \QuaternaryInfC{$\Gamma \tack \text{new } A.k(\overline{e}) : \text{Actor}(A)$}
\end{prooftree}

\vspace{1em}
\begin{prooftree}
    \AxiomC{$\Gamma \tack e_p : T\ \kappa$}
    \AxiomC{$\Gamma \tack e_i : \text{Int}$}
    % \AxiomC{$\Gamma_L = 1$}
    \AxiomC{$\Gamma \tack \text{ref\_readable}(\kappa)$}
    \RightLabel{(Pointer access)}
    \TrinaryInfC{$\Gamma \tack e_p[e_i] : T$}
\end{prooftree}

\vspace{1em}

\begin{prooftree}
    \AxiomC{$\Gamma \tack e : N$}
    \AxiomC{$S(N) = \text{struct } \{ \overline{f:B} \}$}
    \AxiomC{$f_i \in \overline{f}$}
    \RightLabel{(Field access)}
    \TrinaryInfC{$\Gamma \tack e.f_i : B_i$}
\end{prooftree}

\vspace{1em}

\begin{prooftree}
    \AxiomC{$\Gamma \tack e_1 : T_1$}
    \AxiomC{$\Gamma \tack e_2 : T_2$}
    \AxiomC{$\Gamma \tack \text{type\_assignable}(T_1, T_2)$}
    \AxiomC{$\Gamma \tack \text{appear\_lhs}(e_1)$}
    \RightLabel{(Assign)}
    \QuaternaryInfC{$\Gamma \tack e_1 = e_2 : \text{unalias}(T_1)$}
\end{prooftree}

\vspace{1em}

\begin{prooftree}
    \AxiomC{$\Gamma_F(m)_{sig} = (\overline{T_{\text{param}}} \Rightarrow T_{\text{ret}})$}
    \AxiomC{$\Gamma \tack \overline{e} : \overline{T_{\text{arg}}}$}
    \AxiomC{$\Gamma \tack \text{type\_assignable}(\overline{T_{\text{param}}}, \overline{T_{\text{arg}}})$}
    \RightLabel{(Func-Call)}
    \TrinaryInfC{$\Gamma \tack m(\overline{e}) : T_{\text{ret}} $}
\end{prooftree}

\subsection{Statement typing}
$\Gamma \tack s \tackl \Gamma'$ says that $s$ is valid and after that, the environment becomes $\Gamma'$.
A list of statements $\overline{s}$ is typed by the following two rules:

\begin{center}
\begin{tabular}{c@{\qquad\qquad}c}
    \AxiomC{}
    \RightLabel{(Seq-Nil)}
    \UnaryInfC{$\Gamma \tack [] \tackl \Gamma$}
    \DisplayProof
&
    \AxiomC{$\Gamma \tack s \tackl \Gamma'$}
    \AxiomC{$\Gamma' \tack \overline{s} \tackl \Gamma''$}
    \RightLabel{(Seq-Cons)}
    \BinaryInfC{$\Gamma \tack s{::}\overline{s} \tackl \Gamma''$}
    \DisplayProof
\end{tabular}

\vspace{2em}
\end{center}
The typing rules for individual statements are as follows:
\begin{center}
\begin{prooftree}
    \AxiomC{$\Gamma \tack e : T_e$}
    \AxiomC{$\Gamma \tack \text{assignable}(T, T_e)$}
    \RightLabel{(Var-Decl)}
    \BinaryInfC{$\Gamma \tack \text{var } x : T = e; \tackl \Gamma[\Gamma_V \mapsto \Gamma_V[x \mapsto T]]$}
\end{prooftree}

\vspace{2em}

\begin{prooftree}
    \AxiomC{$\Gamma_{\text{this}} = A$}
    \AxiomC{$A_{\text{fields}}(f) = T_f$}
    \AxiomC{$\Gamma \tack e : T_e$}
    \AxiomC{$\text{assignable}(T_f, T_e)$}
    \RightLabel{(Member-Init)}
    \QuaternaryInfC{$\Gamma \tack f := e \tackl \Gamma$}
\end{prooftree}

\begin{prooftree}
    \AxiomC{$\Gamma \tack e_a : \text{Actor}(A_{\text{name}})$}
    \AxiomC{$b \in \Gamma_A(A_{\text{name}})_{behv}$}
    \AxiomC{$\Gamma \tack \overline{e} : \overline{T_e}$}
    \AxiomC{$\text{type\_assignable}(b_{\text{param}}, \overline{T_e})$}
    \RightLabel{[Behv-Call]}
    \QuaternaryInfC{$\Gamma \tack e_a \to b(\overline{e}); \tackl \Gamma$}
\end{prooftree}

\begin{prooftree}
    \AxiomC{$\Gamma \tack e : \text{Int}$}
    \RightLabel{(Out)}
    \UnaryInfC{$\Gamma \tack \text{OUT } e; \tackl \Gamma$}
\end{prooftree}


\begin{prooftree}
    \AxiomC{$\Gamma \tack e : T$}
    \RightLabel{[Raw expression]}
    \UnaryInfC{$\Gamma \tack e; \tackl \Gamma$}
\end{prooftree}


\begin{prooftree}
    \AxiomC{$\Gamma \tack e : \text{bool}$}
    \AxiomC{$\Gamma \tack \overline{s_1} \tackl \Gamma_1$}
    \AxiomC{$\Gamma \tack \overline{s_2} \tackl \Gamma_2$}
    \RightLabel{(If-Else)}
    \TrinaryInfC{$\Gamma \tack \mathtt{if}\ (e)\ \{\ \overline{s_1}\ \}\ \mathtt{else}\ \{\ \overline{s_2}\ \} \tackl \Gamma$}
\end{prooftree}


\begin{prooftree}
    \AxiomC{$\Gamma[L \mapsto \text{true}] \tack \overline{s} \tackl \Gamma'$}
    \RightLabel{[Atomic]}
    \UnaryInfC{$\Gamma \tack \text{atomic } \{ \overline{s} \} \tackl \Gamma$}
\end{prooftree}


\begin{prooftree}
    \AxiomC{$\Gamma \tack e : \text{Bool}$}
    \AxiomC{$\Gamma \tack \overline{s} \tackl \Gamma'$}
    \RightLabel{[While]}
    \BinaryInfC{$\Gamma \tack \text{while } (e) \{ \overline{s} \} \tackl \Gamma$}
\end{prooftree}


\begin{prooftree}
    \AxiomC{$\Gamma \tack e : T$}
    \RightLabel{[Return]}
    \UnaryInfC{$\Gamma \tack \text{return } e; \tackl \Gamma$}
\end{prooftree}
\end{center}

\subsection{Top-level typing}

\begin{center}
\begin{prooftree}
    \AxiomC{$\Gamma \tack \hat{T} : \text{valid}$}
    \RightLabel{[TypeDef]}
    \UnaryInfC{$\Gamma \tack \text{type } N = \hat{T} \tack \Gamma[C \mapsto C[N \mapsto \hat{T}]]$}
\end{prooftree}

\vspace{2em}

\begin{prooftree}
    \AxiomC{$\Gamma' = \Gamma[V \mapsto \{\overline{x : \phi_p}\}]$}
    \AxiomC{$\Gamma' \tack \text{returns}(\overline{s}, \phi_r)$}
    \AxiomC{$\Gamma' \tack \overline{s} \tackl \Gamma''$}
    \RightLabel{[T-Func]}
    \TrinaryInfC{$\Gamma \tack \text{func } m(\overline{x : \phi_p}) \Rightarrow \phi_r \{ \overline{s} \} \tackl \Gamma$}
\end{prooftree}

\vspace{2em}

\begin{prooftree}
    \AxiomC{$\Gamma' = \Gamma[V \mapsto \{\overline{x : \phi_p}\}]$}
    \AxiomC{$\text{doesnt\_return}(\overline{s})$}
    \AxiomC{$\Gamma' \tack \overline{s} \tackl \Gamma''$}
    \RightLabel{[T-Constructor]}
    \TrinaryInfC{$\Gamma \tack \mathtt{new}\ k(\overline{x : \phi_p})\ \{\ \overline{s}\ \} \tackl \Gamma$}
\end{prooftree}

\vspace{2em}

\begin{prooftree}
    \AxiomC{$\Gamma' = \Gamma[V \mapsto \{\overline{x : \phi_p}\}]$}
    \AxiomC{$\text{doesnt\_return}(\overline{s})$}
    \AxiomC{$\Gamma \tack \text{sendable}(\overline{\phi_p})$}
    \AxiomC{$\Gamma' \tack \overline{s} \tackl \Gamma''$}
    \RightLabel{[T-Be]}
    \QuaternaryInfC{$\Gamma \tack \mathtt{be}\ b(\overline{x : \phi_p})\ \{\ \overline{s}\ \} \tackl \Gamma$}
\end{prooftree}

\vspace{2em}

\begin{prooftree}
    \AxiomC{$\Gamma \tack \overline{T} : \text{valid}$}
    \AxiomC{$\Gamma_{\text{base}} = \Gamma[this \mapsto \text{Actor}(A)]$}
    \AxiomC{$\Gamma_{\text{base}} \tack \overline{M} \tackl \Gamma_{\text{base}}$}
    \RightLabel{[T-Actor]}
    \TrinaryInfC{$\Gamma \tack \mathtt{actor}\ A\ \{\ \overline{f : T};\ \overline{M}\ \} \tackl \Gamma$}
\end{prooftree}
\end{center}

\subsubsection{Full program typing}
A full program $P = \overline{D}$ is typed in two phases. First, we collect all top-level declarations of functions and actors to build an initial environment $\Gamma_0$.
\\
In the second phase, the entire declaration list is type-checked under $\Gamma_0$. The program is well-typed if both phases succeed:

\begin{center}
\begin{prooftree}
    \AxiomC{$\text{initial\_env}(\overline{D}) = \Gamma_0$}
    \AxiomC{$\Gamma_0 \tack \overline{D} \tackl \Gamma'$}
    \RightLabel{[T-Program]}
    \BinaryInfC{$\tack \overline{D} : \text{ok}$}
\end{prooftree}
\end{center}

\end{document}
